\section{Microarchitecture}

\begin{definition}
\textbf{Microarchitecture} implements an ISA. One ISA can have many microarchitectures (single-cycle, multi-cycle, pipelined) optimized for \textbf{Performance}, \textbf{Cost} (die area), \textbf{Power/Energy}, and \textbf{Reliability}.
\end{definition}

\subsection{Single-Cycle Microarchitecture}

\begin{definition}
Every instruction takes exactly one clock cycle.
\begin{itemize}
    \item \textbf{CPI}: Always 1.
    \item \textbf{Clock Cycle Time ($T_c$)}: Bounded by the slowest instruction (usually \texttt{lw}).
\end{itemize}
\end{definition}

\subsection{Multi-Cycle Microarchitecture}

\begin{definition}
Breaks instruction execution into multiple shorter cycles. Cycles take variable steps.
Allows \textbf{hardware reuse} (e.g. ALU) and requires internal microarchitectural registers (IR, MDR, ALUOut).
\end{definition}

\begin{example}
Phases: 1. Fetch, 2. Decode, 3. Execute, 4. Memory, 5. Write-back.
\end{example}

\begin{theorem}
$$Time = \#instr \times Average\_CPI \times T_c$$
Multi-cycle has higher CPI but shorter $T_c$, improving performance over single-cycle.
Downside: \textbf{Sequencing overhead} added at each cycle transition.
\end{theorem}

\begin{remark}
Multi-cycle control unit is implemented as an FSM to track instruction steps.
\end{remark}
